% !TeX root = ../tcc.tex
\begin{resumo}
	A hanseníase é uma doença tropical negligenciada cujo diagnóstico se
	baseia principalmente na avaliação clínica, e a semelhança de suas lesões
	com outras condições dermatológicas dificulta a detecção precoce. Modelos de
	visão computacional têm sido propostos para apoiar a identificação dessas
	lesões, mas a forma como seu desempenho é avaliado ainda é pouco examinada.
	Este trabalho analisa a avaliação de desempenho de modelos de visão
	computacional aplicados à identificação de padrões associados à hanseníase
	em imagens de pele e, com base nas evidências reunidas, planeja um estudo de
	ciência de dados. Foi conduzida uma revisão estruturada da literatura,
	adaptada do protocolo de revisão sistemática, com busca assistida por
	ferramentas de inteligência artificial, que incluiu 26 estudos, dos quais 17
	apresentam dados próprios. Predominam as redes convolucionais treinadas por
	transferência e a avaliação em uma única divisão da própria base: dos 16
	estudos que desenvolvem e avaliam modelos, 14 usam essa divisão, apenas 2
	declaram separar as imagens por paciente e só 1 testa o modelo em outra
	base. A acurácia e a sensibilidade são as métricas mais relatadas, nenhum
	estudo reporta o desempenho por tom de pele, e as diferenças entre bases,
	classes e procedimentos impedem comparar numericamente os resultados. Com
	base nessas lacunas, foi planejado, segundo o padrão \textit{Data Science}
	da ACM SIGSOFT, um estudo que avaliará os modelos de fundação DINOv3 e
	PanDerm na distinção entre hanseníase e outras condições de pele em
	fotografias clínicas do \textit{31 Classes of Skin Disease} (31 CSD) e do \textit{Our Normal Skin}, ao lado de um
	classificador trivial e da reprodução de uma abordagem da literatura, com
	partições por grupos de fotografias relacionadas, seleção por validação
	cruzada, acurácia balanceada como medida principal e intervalos de
	confiança por \textit{bootstrap}. O estudo será
	executado na segunda etapa da monografia.

	\vspace{\onelineskip}

	\noindent
	\textbf{Palavras-chave}: hanseníase. visão computacional. avaliação de
	desempenho. aprendizado profundo. revisão estruturada da literatura.
	ciência de dados.
\end{resumo}
